\NSPage{121}%
\subsection{Energy and comparison.}\label{sec:10.3}
In this subsection we first prove that the smooth compactly supported force
\NSi{NS-F-P121-001}{f} of Lemma~\ref{lem:10.3} gives a uniform energy bound for the
constructed velocity \NSi{NS-F-P121-002}{u} up to time one (Lemma~\ref{lem:10.4}).
We then show that any smooth solution with the same force, zero datum, and bounded kinetic
energy agrees with it on every shorter time interval (Lemma~\ref{lem:10.5}). This comparison
will transfer the local velocity growth \eqref{eq:3.6} to a hypothetical global smooth
solution. In these whole-space energy estimates, \NSi{NS-F-P121-003}{\|\cdot\|_p}
denotes the spatial \NSi{NS-F-P121-004}{L^p(\R^3)} norm, and unmarked spatial integrals
are over \NSi{NS-F-P121-005}{\R^3}.

\begin{lemma}\label{lem:10.4}
Let \NSi{NS-F-P121-006}{\mathcal F(t)=\int_0^t\|f(s)\|_2\,ds} for
\NSi{NS-F-P121-007}{0\leq t\leq 1}. Then
\NSi{NS-F-P121-008}{\mathcal F(1)<\infty}, and
\NSFormula{NS-F-P121-009}{display}{10.13}
\begin{equation}
  \|u(t)\|_2^2+2\int_0^t\|\nabla u(s)\|_2^2\,ds
  \leq \mathcal F(t)^2
  \qquad (0\leq t<1).
\tag{10.13}\label{eq:10.13}
\end{equation}
In particular the kinetic energy is uniformly bounded, and the total dissipation on
\NSi{NS-F-P121-010}{[0,1)} is finite.
\end{lemma}

\begin{proof}
On each \NSi{NS-F-P121-011}{[0,T]}, \NSi{NS-F-P121-012}{T<1}, the localized velocity
and pressure are smooth with fixed compact spatial support. Integrating the equation against
\NSi{NS-F-P121-013}{u} therefore gives
\NSFormula{NS-F-P121-014}{display}{10.14}
\begin{equation}
  \frac12\frac{d}{dt}\|u(t)\|_2^2+\|\nabla u(t)\|_2^2
  =\langle f(t),u(t)\rangle.
\tag{10.14}\label{eq:10.14}
\end{equation}
For \NSi{NS-F-P121-015}{\delta>0}, we divide this identity by
\NSi{NS-F-P121-016}{(\|u(t)\|_2^2+\delta^2)^{1/2}}, discard dissipation, and use
Cauchy--Schwarz. Integration from the zero datum and
\NSi{NS-F-P121-017}{\delta\downarrow0} yield
\NSi{NS-F-P121-018}{\|u(t)\|_2\leq\mathcal F(t)}. Returning to
\eqref{eq:10.14} and integrating gives
\NSFormula{NS-F-P121-019}{display}{none}
\[
  \|u(t)\|_2^2+2\int_0^t\|\nabla u(s)\|_2^2\,ds
  \leq 2\int_0^t\mathcal F'(s)\mathcal F(s)\,ds
  =\mathcal F(t)^2.
\]
The compact smooth force has \NSi{NS-F-P121-020}{\mathcal F(1)<\infty}. Monotone
convergence gives finite integrated dissipation on \NSi{NS-F-P121-021}{[0,1)}, using
the bounds for \NSi{NS-F-P121-022}{t<1} alone.
\end{proof}

The comparison argument uses smoothness and a uniform spatial
\NSi{NS-F-P121-023}{L^2} bound on \NSi{NS-F-P121-024}{[0,T]}, where
\NSi{NS-F-P121-025}{T<1}. These hypotheses leave the growth of spatial derivatives at
infinity unrestricted. We recover the pressure gradient from the equation before passing to
the limit in the localized energy identity.

\begin{lemma}\label{lem:10.5}
Fix \NSi{NS-F-P121-026}{T<1}. If \NSi{NS-F-P121-027}{v,P} is a smooth solution of
\eqref{eq:1.1} at viscosity one on \NSi{NS-F-P121-028}{\R^3\times[0,T]}, with the force
\NSi{NS-F-P121-029}{f} of Lemma~\ref{lem:10.3}, zero initial velocity, and
\NSi{NS-F-P121-030}{v\in L^\infty([0,T];L^2(\R^3))}, then
\NSi{NS-F-P121-031}{v=u} on that interval, where \NSi{NS-F-P121-032}{u} is the
localized velocity of Proposition~\ref{prop:10.1}.
\end{lemma}

\begin{proof}
We estimate the difference of the two solutions on expanding balls. The pressure term
requires separate control because no spatial growth condition has been imposed on
\NSi{NS-F-P121-033}{P}. Constants denoted by \NSi{NS-F-P121-034}{C_T} below may
depend on \NSi{NS-F-P121-035}{T} and the stated bounds for the fields, but not on the
cutoff radius. We set
\NSFormula{NS-F-P121-036}{display}{none}
\[
  w=v-u,\qquad
  \pi=P-p,\qquad
  g_{ij}=v_iv_j-u_iu_j=w_iw_j+w_iu_j+u_iw_j.
\]
We use the convention
\NSi{NS-F-P121-037}{(\operatorname{div}g)_i=\sum_j\partial_jg_{ij}}.
On \NSi{NS-F-P121-038}{[0,T]},
\NSi{NS-F-P121-039}{\|w(t)\|_2\leq C_T} and
\NSi{NS-F-P121-040}{\sum_{i,j}\|g_{ij}(t)\|_1\leq C_T}.
The common force cancels, so
\NSFormula{NS-F-P121-041}{display}{10.15}
\begin{equation}
  \partial_t w+(v\cdot\nabla)w+(w\cdot\nabla)u
  =\Delta w-\nabla\pi,\qquad \operatorname{div}w=0.
\tag{10.15}\label{eq:10.15}
\end{equation}

\paragraph{\normalfont\itshape Pressure gradient.}
Let \NSi{NS-F-P121-042}{R_i} be the Riesz transform with Fourier multiplier
\NSi{NS-F-P121-043}{i\xi_i/|\xi|}, and define the distribution
\NSFormula{NS-F-P121-044}{display}{10.16}
\begin{equation}
  \pi_*=\sum_{i,j=1}^3 R_iR_jg_{ij}.
\tag{10.16}\label{eq:10.16}
\end{equation}

\NSPage{122}%
Its multiplier is \NSi{NS-F-P122-001}{-\xi_i\xi_j/|\xi|^2}, with the value at zero
irrelevant. Since the Fourier transform of an \NSi{NS-F-P122-002}{L^1} function is
bounded, \NSi{NS-F-P122-003}{\pi_*} lies uniformly in
\NSi{NS-F-P122-004}{H^{-s}(\R^3)} for each \NSi{NS-F-P122-005}{s>3/2}: its squared
norm is bounded by a constant times
\NSi{NS-F-P122-006}{\|g\|_1^2\int_{\R^3}(1+|\xi|^2)^{-s}\,d\xi}. The multiplier gives
\NSFormula{NS-F-P122-007}{display}{none}
\[
  \Delta\pi_*=-\sum_{i,j}\partial_i\partial_jg_{ij}.
\]
We claim that \NSi{NS-F-P122-008}{\nabla\pi=\nabla\pi_*} in the space--time interior.
For \NSi{NS-F-P122-009}{a\in C_c^\infty(0,T)}, the conservative form of
\eqref{eq:10.15} gives
\NSFormula{NS-F-P122-010}{display}{none}
\[
  \int_0^T a\nabla\pi\,dt
  =\Delta\int_0^T aw\,dt+\int_0^T a'w\,dt
   -\operatorname{div}\int_0^T ag\,dt.
\]
The three terms on the right belong respectively to
\NSi{NS-F-P122-011}{H^{-2}}, \NSi{NS-F-P122-012}{L^2}, and
\NSi{NS-F-P122-013}{H^{-3}}: the last assertion follows from
\NSi{NS-F-P122-014}{g\in L_t^\infty L_x^1} and the Fourier estimate above with
\NSi{NS-F-P122-015}{s=2}. Also
\NSi{NS-F-P122-016}{\pi_*\in L_t^\infty H_x^{-2}}, so
\NSFormula{NS-F-P122-017}{display}{none}
\[
  H_a:=\int_0^T a(\nabla\pi-\nabla\pi_*)\,dt\in H^{-3}(\R^3).
\]
Taking divergence of the difference equation gives
\NSi{NS-F-P122-018}{\Delta\pi=-\sum_{i,j}\partial_i\partial_jg_{ij}=\Delta\pi_*}.
Hence \NSi{NS-F-P122-019}{\Delta H_a=0}. The Fourier transform of
\NSi{NS-F-P122-020}{H_a} is a weighted \NSi{NS-F-P122-021}{L^2} function supported at
\NSi{NS-F-P122-022}{\{0\}}, and is therefore zero. Testing also in space proves the
claim. This determines the pressure gradient under the stated hypotheses, with the spatial
growth of \NSi{NS-F-P122-023}{P} unrestricted.

\paragraph{\normalfont\itshape Pressure flux.}
We now bound the pressure contribution at the boundary of an expanding ball in terms of the
weighted dissipation inside it. We first estimate
\NSi{NS-F-P122-024}{\int\pi w\cdot\nabla\chi_R} by powers of a weighted
\NSi{NS-F-P122-025}{L^6} norm of \NSi{NS-F-P122-026}{w} that can be absorbed by
dissipation. We choose \NSi{NS-F-P122-027}{0\leq\varphi\leq1} smooth, compactly
supported and equal to one on the unit ball. For \NSi{NS-F-P122-028}{R\geq1}, we write
\NSi{NS-F-P122-029}{\varphi_R(x)=\varphi(x/R)}, \NSi{NS-F-P122-030}{\chi_R=\varphi_R^8},
and set
\NSFormula{NS-F-P122-031}{display}{none}
\[
  E_R=\int\chi_R|w|^2,\qquad
  A_R=\left(\int\chi_R|\nabla w|^2\right)^{1/2},\qquad
  B_R=\|\varphi_R^4w\|_6.
\]
The Sobolev inequality and the product rule imply
\NSFormula{NS-F-P122-032}{display}{10.17}
\begin{equation}
  B_R\leq C\|\nabla(\varphi_R^4w)\|_2
  \leq C(A_R+R^{-1}\|w\|_2).
\tag{10.17}\label{eq:10.17}
\end{equation}
We use the standard \NSi{NS-F-P122-033}{L^p} boundedness of Riesz transforms for
\NSi{NS-F-P122-034}{1<p<\infty}; see \cite{ref-20}. For multiplication by a scalar
function \NSi{NS-F-P122-035}{a} and a linear operator \NSi{NS-F-P122-036}{\mathcal T}, we write
\NSi{NS-F-P122-037}{[a,\mathcal T]g=a(\mathcal Tg)-\mathcal T(ag)}. Multiplication of \eqref{eq:10.16} by
\NSi{NS-F-P122-038}{\varphi_R^4} then gives
\NSFormula{NS-F-P122-039}{display}{10.18}
\begin{equation}
  \varphi_R^4\pi_*
  =\sum_{i,j}R_iR_j(\varphi_R^4g_{ij})
   +\sum_{i,j}[\varphi_R^4,R_iR_j]g_{ij}.
\tag{10.18}\label{eq:10.18}
\end{equation}
The first sum has \NSi{NS-F-P122-040}{L^{3/2}} norm at most
\NSi{NS-F-P122-041}{C_T(B_R+1)}, since
\NSFormula{NS-F-P122-042}{display}{none}
\[
  \|\varphi_R^4w_iw_j\|_{3/2}\leq B_R\|w\|_2,\qquad
  \|\varphi_R^4w_iu_j\|_{3/2}\leq\|w\|_2\|u\|_6.
\]
The nonlocal kernel of \NSi{NS-F-P122-043}{R_iR_j} is bounded by
\NSi{NS-F-P122-044}{C|x-y|^{-3}}. The possible multiple of the identity cancels in the
commutator, whose kernel is therefore bounded in absolute value by
\NSFormula{NS-F-P122-045}{display}{none}
\[
  K_R(x-y)=C|x-y|^{-3}\min\{|x-y|/R,1\}.
\]
Direct radial integration gives
\NSFormula{NS-F-P122-046}{display}{none}
\[
  \|K_R\|_{4/3}^{4/3}
  \leq C\left[
    R^{-4/3}\int_0^R r^{-2/3}\,dr+\int_R^\infty r^{-2}\,dr
  \right]
  \leq CR^{-1}.
\]

\NSPage{123}%
Young's convolution inequality bounds the second sum in \eqref{eq:10.18} in
\NSi{NS-F-P123-001}{L^{4/3}} by \NSi{NS-F-P123-002}{C_TR^{-3/4}}. These identities
hold first for compact smooth cutoffs of \NSi{NS-F-P123-003}{g}, then for
\NSi{NS-F-P123-004}{g} by convergence in \NSi{NS-F-P123-005}{L^1}, the uniform
commutator bound, and convergence of its Riesz transforms in
\NSi{NS-F-P123-006}{H^{-s}}. In particular \NSi{NS-F-P123-007}{\pi_*} is locally
integrable in space and time.

The equality of pressure gradients, tested against the compact vector field
\NSi{NS-F-P123-008}{a(t)\chi_Rw}, now gives
\NSFormula{NS-F-P123-009}{display}{none}
\[
  \int\pi w\cdot\nabla\chi_R\,dx
  =\int\pi_*w\cdot\nabla\chi_R\,dx
  \qquad\text{for almost every }t\in(0,T).
\]
Since \NSi{NS-F-P123-010}{|\nabla\chi_R|\leq CR^{-1}\varphi_R^7}, the two terms in
\eqref{eq:10.18} pair with \NSi{NS-F-P123-011}{CR^{-1}\varphi_R^3|w|}.
Interpolation yields
\NSFormula{NS-F-P123-012}{display}{none}
\[
  \|\varphi_R^3w\|_3\leq\|\varphi_R^2w\|_3
  \leq B_R^{1/2}\|w\|_2^{1/2},
  \qquad
  \|\varphi_R^3w\|_4\leq B_R^{3/4}\|w\|_2^{1/4}.
\]
Consequently
\NSFormula{NS-F-P123-013}{display}{10.19}
\begin{equation}
  \left|\int\pi w\cdot\nabla\chi_R\right|
  \leq C_TR^{-1}\left[(B_R+1)B_R^{1/2}+R^{-3/4}B_R^{3/4}\right].
\tag{10.19}\label{eq:10.19}
\end{equation}

\paragraph{\normalfont\itshape Difference energy.}
The pressure flux can now be included in the energy estimate for the difference. All
integrations have compact support, so smoothness suffices. Pairing \eqref{eq:10.15} with
\NSi{NS-F-P123-014}{\chi_Rw} gives
\NSFormula{NS-F-P123-015}{display}{none}
\[
\begin{aligned}
  \frac12E_R'+A_R^2
  ={}&-\int\chi_R(w\cdot\nabla)u\cdot w
      +\frac12\int|w|^2\Delta\chi_R\\
     &+\frac12\int|w|^2v\cdot\nabla\chi_R
      +\int\pi w\cdot\nabla\chi_R.
\end{aligned}
\]
We take \NSi{NS-F-P123-016}{R} large enough that
\NSi{NS-F-P123-017}{\chi_R=1} on a neighborhood of
\NSi{NS-F-P123-018}{\operatorname{supp}u}. Then
\NSi{NS-F-P123-019}{v=w} on \NSi{NS-F-P123-020}{\operatorname{supp}\nabla\chi_R},
and the transport flux is bounded by
\NSFormula{NS-F-P123-021}{display}{none}
\[
  CR^{-1}\int\varphi_R^6|w|^3\leq C_TR^{-1}B_R^{3/2}.
\]
Here
\NSi{NS-F-P123-022}{\varphi_R^2|w|=(\varphi_R^4|w|)^{1/2}|w|^{1/2}} and
Hölder's inequality give the last bound. The Laplacian term is at most
\NSi{NS-F-P123-023}{C_TR^{-2}}. Using \eqref{eq:10.17} in \eqref{eq:10.19}, every
power of \NSi{NS-F-P123-024}{A_R} in these flux bounds is at most
\NSi{NS-F-P123-025}{3/2<2}. Young's inequality therefore gives
\NSFormula{NS-F-P123-026}{display}{none}
\[
  \frac12E_R'+\frac12A_R^2
  \leq\|\nabla u\|_\infty E_R+\frac{C_T}{R}
  \qquad\text{for almost every }t\in(0,T).
\]
For example,
\NSi{NS-F-P123-027}{R^{-1}A_R^{3/2}\leq\delta A_R^2+C_\delta R^{-4}}; the terms
of powers \NSi{NS-F-P123-028}{1/2} and \NSi{NS-F-P123-029}{3/4} have the same
required bound. The coefficient \NSi{NS-F-P123-030}{\|\nabla u\|_\infty} is bounded
on \NSi{NS-F-P123-031}{[0,T]}, and \NSi{NS-F-P123-032}{E_R(0)=0}. Gronwall gives
\NSFormula{NS-F-P123-033}{display}{none}
\[
  E_R(t)\leq\frac{2C_T}{R}\int_0^t
  \exp\left(2\int_s^t\|\nabla u(a)\|_\infty\,da\right)\,ds
  \leq\frac{C_T'}{R}.
\]
Since \NSi{NS-F-P123-034}{\chi_R=1} on each fixed ball for all sufficiently large
\NSi{NS-F-P123-035}{R}, letting \NSi{NS-F-P123-036}{R\to\infty} gives
\NSi{NS-F-P123-037}{w=0} throughout \NSi{NS-F-P123-038}{[0,T]}.
\end{proof}

\subsection{Growth, lifespan, and viscosity.}\label{sec:10.4}
The force extension and energy estimate are now established for the localized fields
\NSi{NS-F-P123-039}{(u,p)} (Lemma~\ref{lem:10.3}, \eqref{eq:10.13}). In this subsection
we use the inner growth asymptotic \eqref{eq:3.6} from Theorem~\ref{thm:3.1} to identify
the maximal classical lifespan and, with the comparison lemma (Lemma~\ref{lem:10.5}),
exclude a global smooth solution of bounded energy. A spatial rescaling then gives the
result for every fixed positive viscosity without changing the singular time.

\NSPage{124}%
\begin{proof}[Proof of Theorem~\ref{thm:1.1}]
The field constructed above solves the equation exactly on
\NSi{NS-F-P124-001}{[0,1)}, by \eqref{eq:10.5}. Smoothness and fixed compact support
give \NSi{NS-F-P124-002}{u\in C([0,T];H^3)} for each
\NSi{NS-F-P124-003}{T<1}. The force has the required smoothness and support by
Lemma~\ref{lem:10.3}, and Lemma~\ref{lem:10.4} gives bounded energy.

For the fixed \NSi{NS-F-P124-004}{X_{\mathrm{in}}} in Theorem~\ref{thm:3.1}, we take
the Cartesian points
\NSFormula{NS-F-P124-005}{display}{10.20}
\begin{equation}
  x_\tau=(\sqrt2X_{\mathrm{in}}\tau,0,0),\qquad t=1-\tau.
\tag{10.20}\label{eq:10.20}
\end{equation}
Along this path \NSi{NS-F-P124-006}{z=0}, \NSi{NS-F-P124-007}{q=\tau}, and the
spatial and time cutoffs in \eqref{eq:10.4} equal one for all sufficiently small
\NSi{NS-F-P124-008}{\tau>0}. Hence \eqref{eq:3.6} gives
\NSFormula{NS-F-P124-009}{display}{10.21}
\begin{equation}
  u_\theta(x_\tau,1-\tau)
  =\tau^{-A}\bigl(e_0+O(\tau^{2h})\bigr)\longrightarrow+\infty.
\tag{10.21}\label{eq:10.21}
\end{equation}
Here the annular corrections vanish in the fixed inner similarity region, and the cutoff sum
of terms with \NSi{NS-F-P124-010}{n\geq1} in the expansion in powers of
\NSi{NS-F-P124-011}{q^{2h}} satisfies the stated relative-error estimate. The summation
cutoffs remain in this estimate. The points \NSi{NS-F-P124-012}{x_\tau} along which the
velocity diverges converge to the origin. The embedding
\NSi{NS-F-P124-013}{H^3(\R^3)\hookrightarrow L^\infty(\R^3)} excludes a classical
\NSi{NS-F-P124-014}{H^3} continuation through time one. In the classical
\NSi{NS-F-P124-015}{H^3} class, the same difference energy calculation is justified by
Sobolev approximation and \NSi{NS-F-P124-016}{H^3\hookrightarrow W^{1,\infty}};
Gronwall gives uniqueness on every shorter interval. Thus the maximal classical existence
interval is \NSi{NS-F-P124-017}{[0,1)}.

Suppose that a global smooth solution \NSi{NS-F-P124-018}{v,P} with the same data had
uniformly bounded kinetic energy. For every \NSi{NS-F-P124-019}{T<1},
Lemma~\ref{lem:10.5} applies on the closed interval \NSi{NS-F-P124-020}{[0,T]}. It
follows that \NSi{NS-F-P124-021}{v=u} throughout \NSi{NS-F-P124-022}{[0,1)}.
Equation~\eqref{eq:10.21} then contradicts the boundedness of the smooth
\NSi{NS-F-P124-023}{v} on a compact neighborhood of \NSi{NS-F-P124-024}{(0,1)}.
The force is nonzero, since \eqref{eq:10.13} would otherwise give
\NSi{NS-F-P124-025}{u=0}.

Finally, for fixed \NSi{NS-F-P124-026}{\nu>0}, we define
\NSFormula{NS-F-P124-027}{display}{10.22}
\begin{equation}
\begin{aligned}
  u_\nu(x,t)&=\sqrt\nu\,u(x/\sqrt\nu,t),\\
  p_\nu(x,t)&=\nu p(x/\sqrt\nu,t),
  &f_\nu(x,t)&=\sqrt\nu\,f(x/\sqrt\nu,t).
\end{aligned}
\tag{10.22}\label{eq:10.22}
\end{equation}
With \NSi{NS-F-P124-028}{y=x/\sqrt\nu}, the chain rule gives
\NSFormula{NS-F-P124-029}{display}{none}
\[
\begin{aligned}
  \partial_tu_\nu+(u_\nu\cdot\nabla_x)u_\nu-\nu\Delta_xu_\nu+\nabla_xp_\nu
  &=\sqrt\nu\,[\partial_tu+(u\cdot\nabla_y)u-\Delta_yu+\nabla_yp](y,t)
    =f_\nu(x,t),\\
  \nabla_x\cdot u_\nu(x,t)&=(\nabla_y\cdot u)(y,t)=0.
\end{aligned}
\]
The datum remains zero and the common spatial support becomes
\NSi{NS-F-P124-030}{\mathcal K_\nu=\sqrt\nu\mathcal K}. Time is unchanged, and
\NSFormula{NS-F-P124-031}{display}{none}
\[
  \partial_x^\alpha\partial_t^m f_\nu(x,t)
  =\nu^{(1-|\alpha|)/2}
   (\partial_y^\alpha\partial_t^m f)(x/\sqrt\nu,t).
\]
Thus the force remains smooth and compactly supported in
\NSi{NS-F-P124-032}{\mathcal K_\nu\times[0,2]}, and satisfies all required decay bounds for each
fixed \NSi{NS-F-P124-033}{\nu>0}. The energy and dissipation scale by
\NSFormula{NS-F-P124-034}{display}{10.23}
\begin{equation}
  \|u_\nu(t)\|_2^2=\nu^{5/2}\|u(t)\|_2^2,\qquad
  \nu\int_0^1\|\nabla u_\nu\|_2^2\,dt
  =\nu^{5/2}\int_0^1\|\nabla u\|_2^2\,dt.
\tag{10.23}\label{eq:10.23}
\end{equation}
Growth occurs at \NSi{NS-F-P124-035}{\sqrt\nu x_\tau} at the same time one. Any smooth
bounded-energy competitor \NSi{NS-F-P124-036}{v_\nu,P_\nu} for the rescaled data would
give the viscosity-one competitor
\NSFormula{NS-F-P124-037}{display}{none}
\[
  v(y,t)=\nu^{-1/2}v_\nu(\sqrt\nu y,t),\qquad
  P(y,t)=\nu^{-1}P_\nu(\sqrt\nu y,t),\qquad
  \|v(t)\|_2^2=\nu^{-5/2}\|v_\nu(t)\|_2^2.
\]
This competitor has already been excluded. The theorem follows for every fixed positive
viscosity.
\end{proof}

\NSPage{125}%
\subsection{The periodic equations.}\label{sec:10.5}
The fixed compact support of the velocity, pressure, and force
(Proposition~\ref{prop:10.1} and Lemma~\ref{lem:10.3}) also allows a periodic
construction. In this subsection we rescale their supports into the interior of a fundamental
cube, then sum their integer translates. The translates have disjoint supports, so the
momentum equation is preserved, including its nonlinear term. The pressure is periodic as
well, as required in the erratum to the problem statement \cite{ref-13}.

The periodic domain and the interior of its fundamental cube are
\NSFormula{NS-F-P125-001}{display}{none}
\[
  \T^3=\R^3/\Z^3,\qquad Q_0=(-1/2,1/2)^3.
\]
The following corollary gives the periodic breakdown alternative (D) in \cite{ref-13}.

\begin{corollary}\label{cor:10.6}
For every \NSi{NS-F-P125-002}{\nu>0}, there is a smooth force on
\NSi{NS-F-P125-003}{\T^3\times[0,\infty)}, compactly supported in time, for which the
solution \NSi{NS-F-P125-004}{(U,P_{\mathrm{per}})} of the periodic Navier--Stokes
equations with viscosity \NSi{NS-F-P125-005}{\nu} and zero initial velocity is smooth on
\NSi{NS-F-P125-006}{[0,1)} and satisfies
\NSFormula{NS-F-P125-007}{display}{none}
\[
  \limsup_{t\uparrow1}\|U(t)\|_{L^\infty(\T^3)}=\infty.
\]
Its velocity and pressure have support in a fixed compact subset of the interior of
\NSi{NS-F-P125-008}{Q_0} at every time before one. There is no global smooth periodic
solution for the same datum and force.
\end{corollary}

\begin{proof}
At a fixed viscosity \NSi{NS-F-P125-009}{\nu}, we take the whole-space fields and force
already constructed and denote them here by \NSi{NS-F-P125-010}{u,p,f}. Their common
spatial support is contained in \NSi{NS-F-P125-011}{\mathcal K_\nu=\sqrt\nu\mathcal K}, as in the
viscosity rescaling \eqref{eq:10.22}. We choose \NSi{NS-F-P125-012}{\lambda>1} so large
that \NSi{NS-F-P125-013}{\lambda^{-1}\mathcal K_\nu\Subset Q_0}, and put
\NSi{NS-F-P125-014}{t_0=1-\lambda^{-2}}. For \NSi{NS-F-P125-015}{t\geq t_0}, we set
\NSFormula{NS-F-P125-016}{display}{none}
\[
\begin{aligned}
  \widetilde u(x,t)&=\lambda u(\lambda x,\lambda^2(t-t_0)),\\
  \widetilde p(x,t)&=\lambda^2p(\lambda x,\lambda^2(t-t_0)),\\
  \widetilde f(x,t)&=\lambda^3f(\lambda x,\lambda^2(t-t_0)).
\end{aligned}
\]
The first two formulas are used for \NSi{NS-F-P125-017}{t<1}, and the force formula is
used for all \NSi{NS-F-P125-018}{t\geq t_0}. We extend all three by zero for
\NSi{NS-F-P125-019}{t<t_0}. They are smooth across
\NSi{NS-F-P125-020}{t=t_0}, since the original fields and force vanish on an interval
after their initial time. Each term of the momentum equation is
\NSi{NS-F-P125-021}{\lambda^3} times its original value, so the viscosity stays equal
to \NSi{NS-F-P125-022}{\nu}. The original time one is reached precisely when
\NSi{NS-F-P125-023}{\lambda^2(t-t_0)=1}, namely at \NSi{NS-F-P125-024}{t=1}.

The periodic fields are defined by
\NSFormula{NS-F-P125-025}{display}{none}
\[
\begin{aligned}
  U(x,t)&=\sum_{k\in\Z^3}\widetilde u(x+k,t),&&0\leq t<1,\\
  P_{\mathrm{per}}(x,t)&=\sum_{k\in\Z^3}\widetilde p(x+k,t),&&0\leq t<1,\\
  F_{\mathrm{per}}(x,t)&=\sum_{k\in\Z^3}\widetilde f(x+k,t),&&t\geq0.
\end{aligned}
\]
Distinct translates have disjoint support, with a positive gap between them. Derivatives
commute with the locally finite sums, and at every point at most one translate is nonzero.
In particular,
\NSFormula{NS-F-P125-026}{display}{none}
\[
  (U\cdot\nabla)U(x,t)
  =\sum_{k\in\Z^3}(\widetilde u(x+k,t)\cdot\nabla)\widetilde u(x+k,t).
\]
Thus the periodized velocity, pressure and force solve the equation exactly:
\NSFormula{NS-F-P125-027}{display}{none}
\[
  \partial_tU+(U\cdot\nabla)U-\nu\Delta U+\nabla P_{\mathrm{per}}
  =F_{\mathrm{per}},\qquad
  \operatorname{div}U=0,\qquad U(\cdot,0)=0.
\]

\NSPage{126}%
The force is smooth, with time support contained in
\NSi{NS-F-P126-001}{[t_0,1+\lambda^{-2}]}. On the torus this implies every decay estimate
\NSi{NS-F-P126-002}{\|\partial_x^\alpha\partial_t^mF_{\mathrm{per}}(t)\|_\infty
\leq C_{\alpha,m,N}(1+t)^{-N}}. For the path \NSi{NS-F-P126-003}{x_\tau} defined in
\eqref{eq:10.20}, the rescaled growth path is
\NSFormula{NS-F-P126-004}{display}{none}
\[
  \widetilde x_\tau=\lambda^{-1}\sqrt\nu x_\tau,\qquad
  \widetilde t_\tau=1-\lambda^{-2}\tau.
\]
It remains inside \NSi{NS-F-P126-005}{Q_0} near \NSi{NS-F-P126-006}{(0,1)}, where the
zeroth translate is the only nonzero term. By \eqref{eq:10.21} and \eqref{eq:10.22},
\NSFormula{NS-F-P126-007}{display}{none}
\[
  U_\theta(\widetilde x_\tau,\widetilde t_\tau)
  =\lambda\sqrt\nu\,\tau^{-A}\bigl(e_0+O(\tau^{2h})\bigr)\longrightarrow+\infty.
\]

Finally, on each \NSi{NS-F-P126-008}{[0,T]} with
\NSi{NS-F-P126-009}{T<1}, two smooth periodic solutions with the same force and datum
have difference \NSi{NS-F-P126-010}{w} obeying
\NSFormula{NS-F-P126-011}{display}{none}
\[
  \frac12\frac{d}{dt}\|w\|_{L^2(\T^3)}^2
  +\nu\|\nabla w\|_{L^2(\T^3)}^2
  \leq\|\nabla U\|_\infty\|w\|_{L^2(\T^3)}^2.
\]
Periodic integration removes both the transport and pressure terms. Gronwall gives equality
of the solutions through \NSi{NS-F-P126-012}{T}. A global smooth competitor would
therefore coincide with \NSi{NS-F-P126-013}{U} before time one and contradict its velocity
blowup at \NSi{NS-F-P126-014}{t=1}.
\end{proof}

\appendix
\section{Matching radial moments and constructing the heat exterior}\label{sec:A}
In this appendix we construct the outer part of the leading profile
\NSi{NS-F-P126-015}{(U,E,\Pi)} in Theorem~\ref{thm:4.6}. Its pressure integral supplies
the datum \NSi{NS-F-P126-016}{\Pi_0} for the regular axis problem
(Lemma~\ref{lem:A.5}). Its moment conditions, including \eqref{eq:A.8}, ensure that the
radially integrated leading residual stress \NSi{NS-F-P126-017}{\mathcal T_0} vanishes in the
exterior after the heat replacement (Lemma~\ref{lem:A.8}). The construction must satisfy
both requirements while retaining the strict stress cone inequalities on the designated
radial intervals (Propositions~\ref{prop:A.4}, \ref{prop:A.7} and \ref{prop:A.10}).

We begin with finite-dimensional moment corrections (Lemmas~\ref{lem:A.1} and
\ref{lem:A.2}) that will also be used to attach the axis profile
(Proposition~\ref{prop:B.8}) and modify the shear (Proposition~\ref{prop:C.2}). We then
build the outer profile from a prescribed inner reference profile (Section~\ref{sec:A.2}).
Its terminal power law \NSi{NS-F-P126-018}{E_{\mathrm{pow}}} is finally replaced by an
exact heat flow, with compensating moment corrections that preserve the pressure datum and
exterior stress conditions (Lemmas~\ref{lem:A.6} and \ref{lem:A.8} and
Proposition~\ref{prop:A.7}).

\subsection{Adjusting finitely many moments.}\label{sec:A.1}
Changes to a profile on a compact radial interval alter its cumulative moments
\eqref{eq:4.15}. In this subsection we correct these changes by adding fixed smooth bumps
with adjustable coefficients. The following results show that distinct power weights give
independent moment changes (Lemma~\ref{lem:A.1}) and that the resulting quadratic moment
equations can be solved for small discrepancies (Lemma~\ref{lem:A.2}). They will allow
each radial modification to preserve the data prescribed farther out, by Lemma~\ref{lem:4.4}.

\begin{lemma}\label{lem:A.1}
Let \NSi{NS-F-P126-019}{\alpha_1,\ldots,\alpha_m} be distinct real numbers. For each
\NSi{NS-F-P126-020}{j} let \NSi{NS-F-P126-021}{\beta_j} be a nonnegative, nonzero
smooth function supported in the interior of a compact interval
\NSi{NS-F-P126-022}{I_j\subset(0,\infty)}, with
\NSi{NS-F-P126-023}{I_1<\cdots<I_m}. The matrix
\NSFormula{NS-F-P126-024}{display}{none}
\[
  B_{ij}=\int_0^\infty x^{\alpha_i}\beta_j(x)\,dx
\]
is invertible. The assertion also holds for exponential weights in a logarithmic coordinate.
If these data vary smoothly over a compact parameter set and the stated separations persist,
the inverse and each fixed parameter derivative of it are bounded.
\end{lemma}
